\section{Campos de fracciones, contenido y lema de Gauss}\label{sec:anillos-fracciones}

\subsection{Ideales extendidos y polinomios primitivos}\label{tema:17-3-1}\label{tema:17-3-2}
Para un ideal $I$ de $R$, el ideal que genera en $R[x]$ es
\[
IR[x]=\set*{\sum_{i=0}^d a_ix^i\given d\geq0,\ a_i\in I}.
\]
En efecto, un polinomio con coeficientes en $I$ es suma de productos $a_ix^i$. A la inversa, multiplicar un elemento de $I$ por un polinomio da coeficientes en $I$, y las sumas conservan esa condición.

\begin{ejer}[Cocientar los coeficientes]\label{ejer:C17-coeficientes}
Probar el isomorfismo
\[
R[x]/IR[x]\cong(R/I)[x].
\]

\end{ejer}
\begin{sol}
Consideremos
\[
\varphi:R[x]\longrightarrow(R/I)[x],\qquad
\sum_i a_ix^i\longmapsto\sum_i(a_i+I)x^i.
\]
Preserva la suma por coeficientes y el producto porque la reducción de cada coeficiente de convolución satisface
\[
\left(\sum_{i+j=k}a_ib_j\right)+I
=\sum_{i+j=k}(a_i+I)(b_j+I).
\]
También preserva la unidad. Es sobreyectivo: se elige un representante de cada uno de los finitos coeficientes de un polinomio del destino. Su núcleo consiste exactamente en los polinomios cuyos coeficientes pertenecen a $I$, es decir, en $IR[x]$. El primer isomorfismo da el resultado.
\end{sol}
\begin{coro}\label{coro:C17-primo-extendido}
Si $I$ es primo en $R$, entonces $IR[x]$ es primo en $R[x]$.
\end{coro}
\begin{proof}
El cociente $(R/I)[x]$ es un dominio porque $R/I$ lo es. Aplicamos el isomorfismo anterior y el criterio de primalidad mediante cocientes.
\end{proof}

En el resto de esta sección $R$ es conmutativo.
\begin{defi}
Un polinomio $f\in R[x]$ es \emph{muy primitivo} si $f\notin\mathfrak pR[x]$ para todo ideal primo $\mathfrak p$ de $R$. Es \emph{primitivo} si se exige lo mismo sólo para los ideales primos principales.
\end{defi}
\begin{lema}\label{lema:C17-producto-primitivo}
El producto $fg$ es primitivo si y sólo si ambos factores lo son.
\end{lema}
\begin{proof}
Para cada ideal primo principal $\mathfrak p$, el \cref{coro:C17-primo-extendido} muestra que $\mathfrak pR[x]$ es primo. Por tanto,
\[
fg\notin\mathfrak pR[x]
\quad\Longleftrightarrow\quad
f\notin\mathfrak pR[x]\ \text{y}\ g\notin\mathfrak pR[x].
\]
Exigirlo para todos esos ideales equivale a la primitividad de ambos factores.
\end{proof}
\begin{lema}\label{lema:C17-criterios-primitivo}
Para $f=a_0+\cdots+a_dx^d$ se cumplen las afirmaciones siguientes.
\begin{enumerate}[label=\textup{(\roman*)}]
\item $f$ es muy primitivo si y sólo si $(a_0,\ldots,a_d)=R$.
\item Si $R$ es DFU y $f\ne0$, es primitivo si y sólo si un máximo común divisor de sus coeficientes es una unidad.
\end{enumerate}
\end{lema}
\begin{proof}
Un ideal primo contiene $f$ coeficiente a coeficiente exactamente cuando contiene $(a_0,\ldots,a_d)$. Si este último es $R$, ningún primo puede contenerlo. Si es propio, la \cref{prop:C17-maximales} lo incluye en un maximal, que también es primo. Esto prueba (i).

Para (ii), si el máximo común divisor no es unidad, tiene un factor irreducible $q$. En un DFU, $(q)$ es primo y contiene todos los coeficientes, luego $f$ no es primitivo. Recíprocamente, si $f\in(q)R[x]$ para un ideal primo principal propio, su generador divide todos los coeficientes. Ese generador no puede ser cero, pues $f\ne0$, ni unidad, pues el ideal es propio. El máximo común divisor tampoco puede ser unidad.
\end{proof}

\subsection{Separar el contenido y multiplicarlo}\label{tema:17-3-3}
\begin{defi}
Si $R$ es DFU y $f\ne0$, su \emph{contenido} $c_f=\operatorname{cont}(f)$ es un máximo común divisor de sus coeficientes. Sólo está definido salvo unidades, pero el ideal principal $C_f=(c_f)$ está determinado de manera única.
\end{defi}
Así, $f$ es primitivo exactamente cuando $C_f=R$. Los productos de contenidos se expresarán mediante ideales: $(a)(b)=(ab)$.

\begin{ejem}
En $\mathbb Z[x]$, el polinomio $f=12x^2-18x+30$ tiene contenido $6$, de modo que
\[
f=6(2x^2-3x+5).
\]
No es primitivo, pero su factor $2x^2-3x+5$ sí lo es, pues $\mcd(2,-3,5)=1$.
\end{ejem}

\begin{lema}\label{lema:C17-parte-primitiva}
Si $R$ es DFU y $f\ne0$, existe $f_0\in R[x]$ primitivo con $f=c_ff_0$. Si $(f)=(c)(g)$ para $c\in R$ y $g$ primitivo, entonces $(c)=C_f$.
\end{lema}
\begin{proof}
Escribamos $a_i=c_fb_i$ para cada coeficiente y pongamos $f_0=\sum_i b_ix^i$. Si un irreducible $q$ dividiera todos los $b_i$, $c_fq$ dividiría todos los $a_i$; por la propiedad del máximo común divisor, $c_fq\mid c_f$. Cancelar $c_f\ne0$ haría a $q$ una unidad. Por tanto, los $b_i$ tienen máximo común divisor unidad y $f_0$ es primitivo.

Para la segunda parte, $(f)=(cg)$ implica $f=dcg$ y $cg=ef$ para ciertos polinomios $d,e$. Como $f\ne0$, cancelar en $f=def$ da $de=1$. En un anillo de polinomios sobre un dominio, las unidades son constantes: $\deg d+\deg e=0$. Así, $d\in R^\times$. Los coeficientes de $cg$ tienen máximo común divisor asociado a $c$, porque los de $g$ no tienen factor irreducible común. Multiplicar por $d$ no cambia su ideal de contenido. Luego $C_f=(c)$.
\end{proof}
\begin{obs}
El contenido se ha definido para polinomios no nulos. Las fórmulas siguientes se utilizan bajo esa hipótesis.
\end{obs}

\begin{prop}[Lema de Gauss]\label{prop:C17-Gauss}
Para $f,g\ne0$ en $R[x]$, con $R$ DFU, se tiene $C_{fg}=C_fC_g$.
\end{prop}
\begin{proof}
Escribimos $f=c_ff_0$ y $g=c_gg_0$, con $f_0,g_0$ primitivos. Su producto también es primitivo por el \cref{lema:C17-producto-primitivo}. La igualdad
\[
(fg)=(c_fc_g)(f_0g_0)
\]
y el \cref{lema:C17-parte-primitiva} dan $C_{fg}=(c_fc_g)=C_fC_g$.
\end{proof}
\begin{coro}\label{coro:C17-contenido-inclusion}
Si $f,g\ne0$ y $(f)\subseteq(g)$ en $R[x]$, entonces $C_f\subseteq C_g$.
\end{coro}
\begin{proof}
La inclusión dice $f=gh$ para algún $h\ne0$. Por Gauss, $C_f=C_gC_h\subseteq C_g$.
\end{proof}

\subsection{La propiedad universal y las fracciones formales}\label{tema:18-1-1}
Sea $R$ un dominio. Consideremos la categoría cuyos objetos son pares $(i,K)$, con $K$ un campo e $i:R\hookrightarrow K$ un homomorfismo inyectivo. Un morfismo de $(i,K)$ a $(j,L)$ es un homomorfismo de campos $\alpha:K\to L$ tal que $\alpha i=j$:
\[
\begin{tikzcd}[column sep=large,row sep=large]
K\arrow[r,"\alpha"]&L\\
R\arrow[u,hook,"i"]\arrow[ur,hook,"j"']&
\end{tikzcd}
\]
Las identidades y composiciones conservan esa condición, así que forman una categoría.

\begin{ejer}[Inyectividad entre campos]\label{ejer:C18-morfismo-campos}
Comprobar que todo homomorfismo unitario entre campos es inyectivo.
\end{ejer}
\begin{sol}
El núcleo es un ideal del campo de partida. No es todo el campo, porque la unidad se envía a una unidad distinta de cero. Como los únicos ideales de un campo son cero y el anillo total, el núcleo es cero. Por tanto, el homomorfismo es inyectivo.
\end{sol}

\begin{defi}
Un \emph{campo de fracciones} de $R$ es un objeto inicial de esta categoría. Lo denotamos $K(R)$, con su inclusión $i:R\hookrightarrow K(R)$.
\end{defi}
Ser inicial significa que para toda inclusión $j:R\hookrightarrow L$, con $L$ campo, existe un único homomorfismo $\widehat j:K(R)\to L$ con $\widehat ji=j$. La unicidad usual de objetos iniciales determina el campo de fracciones salvo isomorfismo compatible con $R$.

Para construirlo, consideremos los pares $(a,r)$ con $a\in R$, $r\ne0$, y la relación
\[
(a,r)\sim(b,s)\quad\Longleftrightarrow\quad as=br.
\]
\begin{ejer}[La equivalencia de fracciones]\label{ejer:C18-equivalencia}
Verificar que $\sim$ es una relación de equivalencia.
\end{ejer}
\begin{sol}
La reflexividad es $ar=ar$, y la simetría se obtiene invirtiendo la igualdad. Para la transitividad, si $as=br$ y $bt=cs$, entonces
\[
s(at-cr)=ast-crs=brt-crs=r(bt-cs)=0.
\]
Como $s\ne0$ y $R$ es dominio, $at=cr$, es decir, $(a,r)\sim(c,t)$.
\end{sol}
Denotamos $a/r$ la clase del par $(a,r)$ y ponemos
\[
K(R)=\set{a/r\given a,r\in R,\ r\ne0}.
\]
Se trata de clases formales, antes de disponer de una división en $R$. Si $s\ne0$, se cumple $as/(rs)=a/r$, porque $(as)r=a(rs)$.

\subsection{Las operaciones y los axiomas de campo}\label{tema:18-1-2}
Definimos
\[
\frac ar+\frac bs=\frac{as+br}{rs},\qquad
\frac ar\frac bs=\frac{ab}{rs}.
\]
Los denominadores son no nulos porque $R$ es dominio. Veamos que las operaciones no dependen de los representantes. Si $ar'=a'r$ y $bs'=b's$, entonces
\[
\begin{aligned}
(as+br)r's'
&=ar'ss'+bs'rr'\\
&=a'rss'+b'srr'=(a's'+b'r')rs,\\
(ab)r's'&=(ar')(bs')=(a'r)(b's)=a'b'rs.
\end{aligned}
\]
Son exactamente las igualdades de productos cruzados requeridas para la suma y el producto.

\begin{ejer}[La estructura de campo]\label{ejer:C18-campo}
Verificar que estas operaciones hacen de $K(R)$ un campo.
\end{ejer}
\begin{sol}
La suma es conmutativa por $as+br=br+as$. Para tres fracciones, ambos modos de asociar la suma dan
\[
\frac{ast+brt+crs}{rst}.
\]
El neutro aditivo es $0/1$, y el opuesto de $a/r$ es $-a/r$, pues su suma es $0/r^2=0/1$. La multiplicación es conmutativa y asociativa porque ambos modos de multiplicar tres fracciones dan $abc/(rst)$; su unidad es $1/1$.

La distributividad se comprueba con denominador común:
\[
\begin{aligned}
\frac ar\left(\frac bs+\frac ct\right)
&=\frac{a(bt+cs)}{rst}
=\frac{abt+acs}{rst}\\
&=\frac{ab}{rs}+\frac{ac}{rt}.
\end{aligned}
\]
La última igualdad se verifica multiplicando numerador y denominador por los factores necesarios; la suma de la derecha tiene inicialmente denominador $r^2st$ y numerador $r(abt+acs)$, que representa la misma clase.

Finalmente, $a/r=0/1$ si y sólo si $a=0$. Toda fracción no nula tiene inversa $r/a$, ya que $(a/r)(r/a)=ar/(ra)=1/1$. Además, $1/1\ne0/1$ porque $1_R\ne0_R$. Se cumplen todos los axiomas de campo.
\end{sol}
La aplicación
\[
i:R\longrightarrow K(R),\qquad a\longmapsto a/1
\]
preserva suma, producto y unidad por las fórmulas anteriores. Es inyectiva porque $a/1=b/1$ implica $a=b$. Identificaremos $R$ con esta copia de sí mismo.

\subsection{El mediador universal y las funciones racionales}\label{tema:18-1-3}
Dada una inclusión $j:R\hookrightarrow L$ en un campo, la fórmula del posible mediador está forzada por
\[
\frac ar=i(a)i(r)^{-1}:
\qquad \widehat j(a/r)=j(a)j(r)^{-1}.
\]
El inverso existe porque $r\ne0$ y $j$ es inyectivo.

\begin{ejer}[El homomorfismo inducido]\label{ejer:C18-mediador}
Verificar que la fórmula anterior define un homomorfismo unitario $\widehat j:K(R)\to L$.
\end{ejer}
\begin{sol}
Si $a/r=b/s$, entonces $as=br$ y, al aplicar $j$ y dividir por $j(r)j(s)\ne0$, obtenemos $j(a)j(r)^{-1}=j(b)j(s)^{-1}$. Por tanto, está bien definido. Para la suma,
\[
\begin{aligned}
\widehat j\left(\frac ar+\frac bs\right)
&=(j(a)j(s)+j(b)j(r))(j(r)j(s))^{-1}\\
&=j(a)j(r)^{-1}+j(b)j(s)^{-1}.
\end{aligned}
\]
Para el producto, $\widehat j(ab/(rs))=j(a)j(b)j(r)^{-1}j(s)^{-1}$ es el producto de las imágenes. También $\widehat j(1/1)=1_L$ y $\widehat j(a/1)=j(a)$, como se requiere.
\end{sol}
\begin{prop}\label{prop:C18-universal}
El campo construido, con $i$, es un campo de fracciones de $R$.
\end{prop}
\begin{proof}
El ejercicio anterior construye el mediador para cualquier objeto $(j,L)$. Todo homomorfismo que extienda $j$ debe enviar $i(a)i(r)^{-1}$ a $j(a)j(r)^{-1}$, por lo que es único. Así se verifica el diagrama universal:
\[
\begin{tikzcd}[column sep=large,row sep=large]
K(R)\arrow[r,dashed,destacado,"\exists!\,\widehat j"]&L\\
R\arrow[u,hook,"i"]\arrow[ur,hook,"j"']&
\end{tikzcd}
\qedhere
\]
\end{proof}

\begin{ejem}
La construcción para $R=\mathbb Z$ da $K(\mathbb Z)=\mathbb Q$. Si $F$ ya es un campo, la aplicación $a/r\mapsto ar^{-1}$ es inversa de $F\hookrightarrow K(F)$. No se añaden elementos a $\mathbb Q$, $\mathbb R$, $\mathbb C$ o $\mathbb Z/p\mathbb Z$, con $p$ primo, al formar su campo de fracciones.
\end{ejem}
\begin{defi}
El \emph{campo de funciones racionales} con coeficientes en un dominio $R$ es $R(x)=K(R[x])$. Sus elementos son fracciones $p(x)/q(x)$ de polinomios de $R[x]$, con $q\ne0$, identificadas por productos cruzados.
\end{defi}
\begin{obs}[Evaluación de fracciones]
La expresión $p(a)/q(a)$ sólo se evalúa donde $q(a)\ne0$. Si $R$ no es campo, su valor pertenece en general a $K(R)$ y no a $R$; por ejemplo, $1/2\notin\mathbb Z$. Además, una función de evaluación puede no determinar la fracción formal.
\end{obs}

\begin{ejer}[Una función no determina el polinomio]\label{ejer:C18-formal-funcion}
Comprobar con un ejemplo sobre un campo finito que dos polinomios formales distintos pueden inducir la misma función.
\end{ejer}
\begin{sol}
En $\mathbb F_2=\mathbb Z/2\mathbb Z$, los polinomios $x^2-x$ y $0$ son distintos: el primero tiene coeficiente principal $1$. Sin embargo, ambos valen cero en los dos elementos $0,1$. También son elementos distintos de $\mathbb F_2(x)$, porque la inclusión del dominio $\mathbb F_2[x]$ en su campo de fracciones es inyectiva. Sus funciones de evaluación coinciden en todo el campo.
\end{sol}

\subsection{Divisibilidad e irreducibilidad sobre el campo de fracciones}\label{tema:18-2-1}\label{tema:18-2-2}
En adelante $R$ es DFU y $K=K(R)$. Para $f\in R[x]$, distinguimos $(f)=fR[x]$ y $(f)_K=fK[x]$. Conservamos la notación $c_f$ para el contenido no nulo y $C_f=(c_f)$ para su ideal.

\begin{lema}\label{lema:C18-descenso}
Sean $f,g\ne0$ en $R[x]$. Si $C_g\subseteq C_f$ y $(g)_K\subseteq(f)_K$, entonces $(g)\subseteq(f)$.
\end{lema}
\begin{proof}
La segunda hipótesis da $g=fh$ con $h\in K[x]$. Podemos escribir
\[
h=\frac ab h_0,\qquad a,b\in R\setminus\{0\},\quad h_0\in R[x]\text{ primitivo}.
\]
En efecto, el producto de los finitos denominadores de sus coeficientes permite escribir $h=H/b$ con $H\in R[x]$ no nulo; separar el contenido de $H$ da esa expresión.

La igualdad $bg=afh_0$ ocurre ya en $R[x]$. Por Gauss y la primitividad de $h_0$,
\[
(b)C_g=(a)C_f.
\]
Como $C_g\subseteq C_f$, resulta $(ac_f)\subseteq(bc_f)$. Hay, pues, $c\in R$ con $ac_f=bc_fc$. Cancelando $c_f\ne0$, obtenemos $a=bc$. Entonces $h=ch_0\in R[x]$ y $g=fh\in(f)$, como se quería.
\end{proof}
El contenido controla precisamente el posible denominador del cociente: la divisibilidad sobre $K$ por sí sola no lo elimina.

\begin{prop}\label{prop:C18-irred-fracciones}
Sea $f\in R[x]$ de grado positivo. Si es irreducible sobre $R$, entonces es primitivo y también es irreducible sobre el campo de fracciones $K$.
\end{prop}
\begin{proof}
Si $f=c_ff_0$ no fuera primitivo, $c_f$ sería una constante no unidad y $f_0$ seguiría teniendo grado positivo. Ambos serían no unidades de $R[x]$, contradiciendo la irreducibilidad de $f$.

Supongamos ahora $f=gh$ en $K[x]$. Como en la prueba anterior, escribimos $g=cg_0$, $h=dh_0$, con $c,d\in K^\times$ y $g_0,h_0\in R[x]$ primitivos. Su producto es primitivo, y
\[
C_f=C_{g_0h_0}=R,\qquad (f)_K=(g_0h_0)_K.
\]
El \cref{lema:C18-descenso}, aplicado en ambos sentidos, da $(f)=(g_0h_0)$ en $R[x]$. Por tanto, $f=ug_0h_0$ para una unidad $u\in R[x]$. La irreducibilidad obliga a que $g_0$ o $h_0$ sea unidad. Entonces $g$ o $h$ es unidad en $K[x]$, y $f$ es irreducible allí.
\end{proof}
Ser primitivo por sí solo no implica irreducibilidad: $x^2-1=(x-1)(x+1)$ es primitivo y reducible en $\mathbb Z[x]$.

\begin{coro}\label{coro:C18-criterio-irred}
Un polinomio no constante $f\in R[x]$ es irreducible en $R[x]$ si y sólo si es primitivo e irreducible en $K[x]$.
\end{coro}
\begin{proof}
Una dirección es la proposición. Para la otra, sea $f=gh$ con $g,h\in R[x]$. La irreducibilidad sobre $K$ hace que uno de ellos, digamos $g$, sea una constante no nula. Gauss da
\[
R=C_f=(g)C_h\subseteq(g)\subseteq R.
\]
Luego $g$ es unidad de $R$, y también de $R[x]$. No puede haber una factorización en dos no unidades.
\end{proof}

\subsection{Factorialidad de los anillos de polinomios}\label{tema:18-3-1}\label{tema:18-3-2a}
\begin{teo}\label{teo:C18-polinomios-DFU}
Si $R$ es DFU, entonces $R[x]$ es DFU.
\end{teo}
\begin{proof}
Usamos el \cref{teo:C16-criterio}: condición de cadena ascendente para ideales principales y primalidad de todos los irreducibles.

Sea $(f_1)\subseteq(f_2)\subseteq\cdots$ una cadena de principales. Si todos los generadores son cero, ya se estabiliza. En otro caso, descartamos los términos iniciales nulos. Por el \cref{coro:C17-contenido-inclusion}, se obtiene una cadena
\[
C_{f_1}\subseteq C_{f_2}\subseteq\cdots
\]
de ideales principales en $R$, que se estabiliza porque $R$ es DFU. También se estabiliza la cadena $(f_1)_K\subseteq(f_2)_K\subseteq\cdots$ en $K[x]$, que es DIP y por tanto DFU. Para índices suficientemente grandes tenemos ambas igualdades consecutivas. El \cref{lema:C18-descenso} da entonces $(f_{i+1})\subseteq(f_i)$; la inclusión contraria ya pertenecía a la cadena. Así se estabiliza la cadena original.

Sea ahora $f\in R[x]$ irreducible. Si es constante, es irreducible en $R$: una factorización por constantes no unidades también lo sería en $R[x]$. Es primo en $R$ porque $R$ es DFU. El \cref{coro:C17-primo-extendido} muestra que $(f)$ es primo en $R[x]$.

Si $f$ no es constante, es primitivo e irreducible en $K[x]$ por la \cref{prop:C18-irred-fracciones}. Como $K[x]$ es DIP, $(f)_K$ es primo. Consideremos
\[
\rho:R[x]\hookrightarrow K[x]\longrightarrow K[x]/(f)_K.
\]
Su núcleo es $(f)$. Una inclusión es inmediata. Para la otra, si $\rho(g)=0$ y $g\ne0$, entonces $(g)_K\subseteq(f)_K$ y $C_g\subseteq R=C_f$. El \cref{lema:C18-descenso} implica $g\in(f)$; para $g=0$ es inmediato. El primer isomorfismo induce la inclusión
\[
\begin{tikzcd}[column sep=large]
R[x]\arrow[r,hook]\arrow[d,two heads]&K[x]\arrow[d,two heads]\\
R[x]/(f)\arrow[r,hook,destacado]&K[x]/(f)_K.
\end{tikzcd}
\]
El anillo de abajo a la derecha es dominio, luego también lo es el de la izquierda. Por tanto, $(f)$ es primo. El criterio citado concluye que $R[x]$ es DFU.
\end{proof}

\Needspace{6\baselineskip}
\begin{ejer}[Elegir factores sobre el dominio]\label{ejer:C18-elegir-factores}
Explicar en qué sentido los factores irreducibles de un polinomio $f\ne0$ de $R[x]$, calculados en $K[x]$, pueden elegirse en $R[x]$.
\end{ejer}
\begin{sol}
Por el teorema, podemos factorizar $f$ en $R[x]$. Separamos los factores constantes, que se vuelven unidades en $K[x]$, de los factores no constantes $p_1,\ldots,p_s$. Estos últimos son primitivos e irreducibles en $K[x]$ por la \cref{prop:C18-irred-fracciones}. Por tanto,
\[
f=\lambda p_1\cdots p_s,\qquad \lambda\in K^\times,
\]
es una factorización sobre $K$. La unicidad en $K[x]$ dice que cualquier otra lista de factores irreducibles coincide con ésta salvo orden y multiplicación por constantes de $K^\times$. Así, cada factor no constante puede reemplazarse por un asociado que pertenece a $R[x]$, conservando aparte el factor escalar. No se afirma que cada representante inicial tenga ya coeficientes en $R$.
\end{sol}
